This module controls the position and orientation of an actor along a
predefined trajectory. It can operate independently of the session
timeline using an internal clock, or optionally synchronize directly
with the \tascar{} time. Orientation can be explicitly defined or
automatically calculated based on the direction of movement along the
path. The module exposes the following OSC variables (see below for
more details):

\begin{verbatim}
/motionpath/go from to
/motionpath/start
/motionpath/stop
/motionpath/locate time
/motionpath/stoptime time
\end{verbatim}
{\tt go} moves along the path from the time {\tt from} until the time {\tt to}.
%
{\tt start} starts the motion at the current time, {\tt stop} stops the motion.
%
{\tt locate} sets the path time to {\tt time} without changing the motion state.
%
{\tt stoptime} sets the time when the motion will be stopped. If the
current time is after the stop time, then the current time is set to
the stop time.

Additionally, the module provides an {\tt /active} OSC command to
enable or disable the plugin. When disabled, the motion path is
ignored and the actor remains in its current state. The specific OSC
paths are constructed using the identifier defined in the
configuration (defaulting to ``motionpath''), allowing multiple
instances to be controlled independently.

The spatial trajectory is defined using a \elem{position} track (see
Section \ref{sec:position}), while orientation can be specified via an
\elem{orientation} (see Section \ref{sec:orientation})
track. Alternatively, orientation can be derived automatically from
the movement direction using the \attr{sampledorientation}
parameter. When this parameter is non-zero, the actor faces along the
tangent of the path, looking ahead by the specified distance in meters
(or behind if the value is negative). The module also supports a
\elem{creator} XML structure, which generates path points and
automatically calculates orientation based on the direction of
movement between those points, handling azimuth unwrapping to prevent
discontinuities.

\definecolor{shadecolor}{RGB}{255,230,204}\begin{snugshade}
{\footnotesize
\label{attrtab:motionpath}
Attributes of element {\bf motionpath}\nopagebreak

\begin{tabularx}{\textwidth}{l>{\raggedright}XX}
\hline
name & description (type, unit) & def.\\
\hline
\hline
\indattr{active} & Enables or disables the plugin. If set to false, the motion path is ignored and the actor remains in its current state. (bool) & true\\
\hline
\indattr{actor} & pattern to match actor objects (string array) & \\
\hline
\indattr{id} & The identifier used for OSC control paths (e.g., /motionpath/go). (string) & \\
\hline
\indattr{sampledorientation} & Defines how orientation is calculated. If 0, orientation is taken from the <orientation> track. If non-zero, orientation is calculated based on the direction of movement (tangent) looking ahead (positive value) or behind (negative value) by the specified distance along the path in meters. (double, m) & 0\\
\hline
\indattr{tascartime} & Determines the time reference. If true, the path position is synchronized with the main TASCAR session time. If false, the plugin uses its own internal clock, controllable via OSC (/go, /start, /stop). (bool) & false\\
\hline
\end{tabularx}
}
\end{snugshade}


\definecolor{shadecolor}{RGB}{236,236,255}\begin{snugshade}
{\footnotesize
\label{osctab:motionpath}
OSC variables:
\nopagebreak

\begin{tabularx}{\textwidth}{llllX}
\hline
path & fmt. & range & readable & description\\
\hline
\attr{/.../active} & i & bool & yes & \\
\attr{/.../go} & ff &  & no & \\
\attr{/.../locate} & f &  & yes & \\
\attr{/.../start} &  &  & no & \\
\attr{/.../stop} &  &  & no & \\
\attr{/.../stoptime} & f &  & yes & \\
\hline
\end{tabularx}
}
\end{snugshade}
\definecolor{shadecolor}{RGB}{255,230,204}

